\subsection*{Lower image likelihood loss coexists with more empty book outputs}

Allowing the continuation query and residuals to learn through the joint image
likelihood completes the same 13,766-source, 41,298-update schedule and
6,108,741 target exposures. The independent source/optimizer audit, five fixed
snapshots and all 124 conditional cases are complete. Original frozen candidate
credits are 970/1,040, compared with 978/1,040 for the unchanged head-only pixels
reviewed in this comparison. Candidate book groups receive 93/108 development
credits, 95/108 training-fixed credits and 92/108 training-proportional credits;
all miss the unchanged 90\% prerequisite. QA groups pass numerically, at
250/257 and 251/257 for development fixed/proportional inputs, and 94/101 and
95/101 for the corresponding training groups. The actor therefore stops before
autonomous evaluation. Its complete-language quality remains unmeasured.

All 33 boards and 2,080 readings were frozen before owner unmasking. The
subsequent pixel audit detects 28 contradictory blank flags, including four
candidate and two control credited empty images. Original scores and uncertain
readings remain unchanged. These credits are retained review outcomes rather
than corrected physical accuracy. The complete-array measurements below use
actual binary ink and do not repair any transcription.

Source-image likelihood improves while empty output becomes more frequent
(Fig.~\ref{fig:context-likelihood-failure}a,b). Development-book mean NLL per
physical mask decreases from 0.107162 to 0.076070, but false blanks increase
from 42 to 93 among 1,161 positions. Training-fixed NLL decreases from 0.076528
to 0.055277 while false blanks increase from 137 to 307; training-proportional
NLL decreases from 0.076529 to 0.055308 while false blanks increase from 145
to 289. Each training layout retains 4,586 positions. Truly empty sources are
counted separately: 22 development and 110 per training layout. Across book
groups, false blanks rise from 324 to 689; across all seven groups, from 349
to 711. Every third-epoch book presentation also has lower online source NLL
than its matched head-only presentation. This loss improvement does not
qualify generated writing or language.

The completed actual-context probability study restores all 131 learned tensors
and recomputes all 12,597 target-preceding states. Only 57/711 false blanks
prefer their source image to the empty image by the fixed 0.001-nat margin.
Observed book counts are 4/93, 26/307 and 27/289, below the prospective
requirements of 47, 154 and 145 (Fig.~\ref{fig:context-likelihood-failure}c).
A first-nonzero-mask upper bound supports an empty global mode at 108 false
blanks, including 102 book positions. All saved physical argmax choices
reproduce without disagreement; maximum mean source-NLL drift is
$6.222\times10^{-9}$. The bound remains subject to finite-precision arithmetic.
These results reject the registered search-only eligibility rule and motivate
a training-objective comparison; they do not identify a unique cause of failed
autonomous language. The six original examples in
Fig.~\ref{fig:context-likelihood-failure}d follow fixed case/position order.

The context-fit worker takes 76,155.600 seconds, including 1,701.612 seconds of
resource waiting, at 18.281 GiB peak allocated VRAM. Coordinator-measured
training, audit and qualification child times are 76,158.268, 15.255 and
477.432 seconds. Inherited pretraining, preparation and prior trials remain
separate costs. The new whole-glyph ranking objective passes a reset
real-source preflight at 17.795 GiB, but its separately registered full-corpus
comparison has no measured writing or complete-language result here.
